\scanpage{19}
définit $u+v$ $(u,v\in V^*)$ séparément quand $u$ ou $v=0$ (comme étant $v$ resp. $u$), quand $u,v\in V^*$ et $u,v$ dans un $=$ $D_x$ (par la loi additive de $D_x$), et quand $u+v$ sont disjoints par \deleted{la \unsure{pr|p}} comme $-w$, où $w$ est l'\uline{unique} $w\in V^*$ qui soit disj. de $u,v$ et tel que $u\wedge v=v\wedge w=w\wedge u$. Pour pouvoir réaliser cette construction, on a besoin des \deleted{ceci} propriétés que voici sur les données a) b) c) [d)] e) :

a) Pour $x,y,z$ \inserted{$\in P$} mutuellement disjoints, on peut trouver $u\in T_x$, $v\in T_y$, $w\in T_z$ tels que l'on ait
\[
u\wedge v=v\wedge w=w\wedge u\qquad(\text{soit }\alpha\mathrel{\unsure{+|\in}}\omega)
\]
(et alors on peut $=$ supposer que $u$ \inserted{ou} $v$ \inserted{ou} $w$ \deleted{\unsure{ou $\alpha$|ou $x$}} soit donné d'avance).
% archive sheet number: 18
